\begin{musica}{Loucos de Lisboa}{João Gil, João Monge}{Tuna de Engenharia da Universidade do Porto}
\nota{Introdução: [G C G D G A7 C D G] 2x}
\nota{\phantom{Introdução: }C D G C D G C D G C D G D7}

Parava \ac{G}{no} café quando eu lá es\ac{C}{ta}va,
Na voz ti\ac{G}{nha} o talento dos pe\ac{D}{din}tes,
Entre um ci\ac{G}{ga}rro e outro lá cra\ac{A7}{va}va a bica,
\ac{C}{Ao} me\ac{D}{lhor} dos seus ou\ac{G}{vin}tes.

As mãos e o olhar da mesma cor,
Cinzenta como a roupa que trazia;
Num gesto que podia ser de amor
Sorria, e ao sorrir agradecia.

\nota{Refrão:}
São os lou\ac{C}{cos} de Lisboa \ac{D}{}_
Que nos fa\ac{G}{zem} duvidar,
A Terra \ac{C}{gi}ra ao contrário \ac{D}{}_
E os rios \ac{G}{nas}cem no mar.

\nota{Instrumental: C D G C D G G7}
\nota{\phantom{Instrumental: }C D G E7 Am D G D7}

Um dia numa sala do quarteto
Passou um filme lá do hospital,
Onde o esquecido filmado no gueto
Entrava como artista principal.

Comprámos a entrada p'ra sessão
P'ra ver tal personagem no ecrã,
O rosto maltratado era a razão
De ele não aparecer pela manhã.

\nota{Refrão}

\nota{Instrumental: C D G C D G G7}
\nota{\phantom{Instrumental: }C D G E7 Am D G D7}

Mudámos muita vez de calendário
Como o café mudou de freguesia,
Deixámos de tributo a quem lá para
Um louco a fazer-lhe companhia.

É sempre a mesma pose o mesmo olhar
De quem não mede os dias que vagueiam
Sentado lá continua a cravar beijinhos
Às meninas que passeiam.

\nota{Refrão}

\nota{Instrumental: C D G C D G}

\nota{Refrão (só vozes)}

\nota{Refrão}

\posicoes{C: x32010, D: xx0232, D7: xx0212, E7: 020100, G: 320033, Am: x02010}
\end{musica}
